\scanpage{30}
% Pencil archive foliation: 29.
8) Soient $\ell,\ell'\in\mathbb{P}$ premiers aux car rés.
de $X$. Soit $M\in\operatorname{Ob}\mathcal{M}^+(X)$. On veut
que
\[
\begin{array}{c}
T_\ell(M)\ \text{faisceau constant tordu}\\
\Updownarrow\\
T_{\ell'}(M)\ \text{faisceau constant tordu}
\end{array}
\]
et que sous ces conditions, $T_\ell(M)$ et
$T_{\ell'}(M)$ doivent avoir \unsure{$=$|$\overline{\mathrm m}$} rang en chaque
pt. Pour vérifier l'égalité des rangs, on
est ramené au cas où $X$ est \inserted{le} spectre d'un
corps (\unsure{fini} si on veut, \deleted{\unsure{à}} ou clôture alg.
d'un tel). Plus généralement, si
$u$ est un endomorphisme de $M$,
($X=\operatorname{Spec}k$, $k$ un corps), on en déduit
\[
T_\ell(u)\in\operatorname{End}T_\ell(M),\qquad
T_{\ell'}(u)\in\operatorname{End}(T_{\ell'}(M))
\]
et je dis qu'on a
\[
\boxed{\operatorname{Tr}T_\ell(u)=\operatorname{Tr}T_{\unsure{\ell|\ell'}}(u)
\in\mathbb{Q}}
\]
d'où, remplaçant $u$ par $\unsure{\bigwedge^i u|\lambda^i u}$, le fait
\[
\boxed{P(T_\ell(u),t)=P(T_{\ell'}(u),t)\in\mathbb{Q}[t]}
\]
